\paragraph*{Data.} Spin configurations are generated with a batched Swendsen--Wang cluster sampler~\cite{swendsen1987nonuniversal} in NumPy/SciPy (bonds between equal neighbours are activated with probability $1-e^{-2/T}$, clusters found by connected components over the whole batch, each flipped with probability $1/2$). We chose it over Wolff~\cite{wolff1989collective} because one call advances every temperature and chain at once; a checkerboard Metropolis sampler is also implemented and used only as a cross-check. For each lattice size there are \rhval{const/n_temps} temperatures, linearly spaced in $[1.5,3.5]$, and \rhval{const/n_chains} independent chains per temperature, each started from a random configuration, thermalised for \rhval{const/sw_therm} sweeps and sampled \rhval{const/n_per_chain} times at a gap of \rhval{const/sw_gap} sweeps, i.e.\ \rhval{const/n_samples_per_t} samples per temperature. Table~\ref{tab:sampler} compares the mean energy per site of the sampler, of checkerboard Metropolis (cold start) and of Onsager's exact infinite-lattice value on $L=24$ ($64$ chains per temperature for this check); the lag-1 autocorrelation of $|m|$ between stored samples at $L=32$ near $T_c$ is \rhval{sanity/sampler-check-long/ising/lag1_autocorr_absm_l32_tc/mean:2}. A first, smaller check with fewer chains is also in the run registry (\texttt{Sampler check}).

\begin{table}[t]\centering\caption{Sampler check, $L=24$: energy per site from Swendsen--Wang, checkerboard Metropolis, and the exact infinite-lattice value (finite-size and statistical differences are of order $10^{-3}$).}\label{tab:sampler}
\begin{tabular}{lccc}
\toprule
$T$ & $e$ Swendsen--Wang & $e$ Metropolis & $e$ Onsager \\
\midrule
1.8 & \rhval{sanity/sampler-check-long/ising/e_sw_tlong_1p8/mean:4} & \rhval{sanity/sampler-check-long/ising/e_metropolis_tlong_1p8/mean:4} & \rhval{sanity/sampler-check-long/ising/e_onsager_tlong_1p8/mean:4} \\
2.0 & \rhval{sanity/sampler-check-long/ising/e_sw_tlong_2p0/mean:4} & \rhval{sanity/sampler-check-long/ising/e_metropolis_tlong_2p0/mean:4} & \rhval{sanity/sampler-check-long/ising/e_onsager_tlong_2p0/mean:4} \\
3.0 & \rhval{sanity/sampler-check-long/ising/e_sw_tlong_3p0/mean:4} & \rhval{sanity/sampler-check-long/ising/e_metropolis_tlong_3p0/mean:4} & \rhval{sanity/sampler-check-long/ising/e_onsager_tlong_3p0/mean:4} \\
3.5 & \rhval{sanity/sampler-check-long/ising/e_sw_tlong_3p5/mean:4} & \rhval{sanity/sampler-check-long/ising/e_metropolis_tlong_3p5/mean:4} & \rhval{sanity/sampler-check-long/ising/e_onsager_tlong_3p5/mean:4} \\
\bottomrule
\end{tabular}

\end{table}

\paragraph*{Splits and seeds.} Chains 0--2 (\rhval{const/n_per_chain} samples each) train or fit; chain 3 is held out, so every output curve is evaluated on a chain no model has seen, including at the temperatures inside the training windows. Seeds $0$--$9$ regenerate the data (seed $\rhval{const/data_seed_stride}\,\mathrm{seed}+L$) and the model initialisation; the main group has \rhval{count/main/seeds} seeds $\times$ \rhval{count/main/systems} systems.

\paragraph*{Hyperparameters and tuning.} Adam, learning rate \rhval{const/lr:3}, weight decay $10^{-4}$, batch \rhval{const/batch_size}, \rhval{const/epochs} epochs for supervised models; \rhval{const/conf_hidden} hidden units and \rhval{const/conf_epochs} epochs for each confusion network. These were fixed a priori and not tuned on any $T_c$ quantity (a pilot on an unreported seed was used only to check run time and that training converged); the supervised models see no label information from the transition region, and the windows are placed using the rounded prior $T_c^{\rm ref}$. Baselines were therefore not tuned either: that is a limitation (Sec.~\ref{sec:limitations}). Learning-by-confusion networks are smaller than the supervised ones by design, because $\approx$32 of them are trained per size.

\paragraph*{Metrics.} Absolute error $|\hat T_c-T_c|$ per size (\texttt{tc\_error\_l16/24/32}) and after the collapse (\texttt{tc\_error\_fss}); the signed bias; $|\hat\nu-1|$; and for the supervised models the held-out accuracy in the training windows (\texttt{acc\_far}). All systems share the same data, temperature grid, held-out chain and metric definitions. Hardware: CPU only, two threads per process, at most two processes at a time; every run completed well inside the 6-minute limit. Analyses of saved curves (collapse variants) and the paired tests are logged runs too.
